\documentclass{article}
\usepackage{enumitem}
\usepackage{amsmath}
\begin{document}

\title{Chapter 12: Mineral Nutrition}
\date{}
\maketitle

\begin{enumerate}[label=Q\arabic*.]

\item Which of the following criteria must ALL be satisfied for an element to be called ``essential'' for plants?
\\(A) Present in large amounts and improves growth
\\(B) Deficiency causes abnormal growth; element must be part of an essential molecule or have a specific metabolic role; no other element can substitute
\\(C) Found in all plant tissues and improves yield
\\(D) Present in soil and absorbed through roots

\item Deficiency of which mineral causes interveinal chlorosis starting in older (mature) leaves due to its mobility in the plant?
\\(A) Iron
\\(B) Calcium
\\(C) Magnesium
\\(D) Sulfur

\item Nitrogen fixation in legumes is carried out by \textit{Rhizobium} in root nodules. The enzyme responsible is:
\\(A) Nitrate reductase
\\(B) Nitrogenase
\\(C) Glutamine synthetase
\\(D) Amylase

\item The ``leghemoglobin'' in root nodules functions to:
\\(A) Transport nitrogen gas to nitrogenase
\\(B) Scavenge oxygen to maintain the anaerobic conditions required by nitrogenase
\\(C) Provide carbon skeletons for amino acid synthesis
\\(D) Transport fixed nitrogen to leaves

\item Assertion: Deficiency of calcium causes stunted growth of meristematic tissue.\\
Reason: Calcium is an immobile element needed for middle lamella formation (calcium pectate) and cell division.
\\(A) Both true, Reason is correct explanation
\\(B) Both true, Reason is not correct explanation
\\(C) Assertion true, Reason false
\\(D) Assertion false, Reason true

\item Which of the following elements is needed in trace (micro) amounts but is essential for nitrogen fixation?
\\(A) Nitrogen
\\(B) Molybdenum
\\(C) Phosphorus
\\(D) Potassium

\item Hydroponics is used to study mineral nutrition because:
\\(A) Plants grow faster in water than soil
\\(B) Specific minerals can be omitted to study their deficiency symptoms
\\(C) It requires no energy for nutrient absorption
\\(D) Roots absorb more minerals in solution

\item Which process converts nitrate (NO$_3^-$) back to atmospheric nitrogen (N$_2$) in the nitrogen cycle?
\\(A) Ammonification
\\(B) Nitrification
\\(C) Denitrification
\\(D) Nitrogen fixation

\item Deficiency of phosphorus leads to:
\\(A) Yellowing of young leaves
\\(B) Stunted root and shoot growth, purpling of leaves due to anthocyanin accumulation
\\(C) Tip burn of leaves
\\(D) Interveinal chlorosis in old leaves

\item \textit{Nepenthes} (pitcher plant) is an example of an insectivorous plant that:
\\(A) Fixes nitrogen using symbiotic bacteria
\\(B) Supplements nitrogen nutrition by digesting trapped insects
\\(C) Absorbs mineral nitrogen from soil exclusively
\\(D) Synthesizes nitrogen from atmospheric CO$_2$

\item The process of nitrification involves conversion of:
\\(A) N$_2$ to NH$_3$ by \textit{Rhizobium}
\\(B) NH$_4^+$ to NO$_2^-$ to NO$_3^-$ by \textit{Nitrosomonas} and \textit{Nitrobacter}
\\(C) NO$_3^-$ to N$_2$ by \textit{Pseudomonas}
\\(D) Organic nitrogen to NH$_3$ by decomposers

\item Which macronutrient is a component of ATP, NADP, and nucleic acids?
\\(A) Nitrogen
\\(B) Sulfur
\\(C) Phosphorus
\\(D) Potassium

\item Iron deficiency in plants causes chlorosis in young leaves (not old leaves) because:
\\(A) Iron is mobile in the phloem
\\(B) Iron is immobile and cannot be remobilized from old leaves to young leaves
\\(C) Iron is destroyed by sunlight in young leaves
\\(D) Young leaves have higher iron demand than old leaves

\item Match the following minerals with their deficiency symptoms:\\
1. Nitrogen $\rightarrow$ (a) Yellowing of older leaves (chlorosis), stunted growth\\
2. Iron $\rightarrow$ (b) Interveinal chlorosis in young leaves\\
3. Calcium $\rightarrow$ (c) Death of growing tips (meristems), stunted roots\\
4. Potassium $\rightarrow$ (d) Scorching/browning of leaf margins
\\(A) 1-a, 2-b, 3-c, 4-d
\\(B) 1-b, 2-a, 3-d, 4-c
\\(C) 1-c, 2-d, 3-a, 4-b
\\(D) 1-d, 2-c, 3-a, 4-b

\item The term ``biofertilizer'' refers to:
\\(A) Chemical fertilizers produced by biotechnology
\\(B) Living organisms (bacteria, cyanobacteria, fungi) that enhance soil nutrient availability
\\(C) Compost made from plant material
\\(D) Slow-release inorganic fertilizers

\end{enumerate}
\end{document}
